\documentclass[11pt]{article}
\usepackage[margin=23mm]{geometry}
\usepackage{amsmath,amssymb,booktabs}
\title{Mixed four-loop two-point forests\\with one weighted bubble edge}
\author{Phi3Lab isolated development}
\date{25 October 2026}
\begin{document}
\maketitle
\noindent Targets: L4E2G6, G19, G15, G21, G46, G63. All expressions use
unit mass and measure $d^6k/\pi^3$ per loop, without coupling or graph
symmetry weights. There are 84 original forests and 42 bubble pairs.
The graph audit checks all $6(2^{11}-1)=12282$ nonempty edge subsets.
These are algebraic identities, not complete-integral precision claims.

\paragraph{Original edge identities and changing bases.}
A vertex subtraction can make a previously selected edge basis singular.
For each compatible vertex forest, choose four original edge routes with
loop coefficient matrix $L_f$ satisfying $\det L_f=\pm1$.
The affine shift includes the retained external momentum. The same basis
is used in the bubble-present and bubble-subtracted member of the pair.
The two propagators carrying the bubble momentum are identified in the
bare graph; additional equal routes after a parent projection are retained
as separate edges. All transformed Gaussian determinants are compared
exactly to the original positive polynomials, including Taylor variables.

\paragraph{Bubble mass and kinetic subtraction.}
For $t=k^2$, the degree-two bubble remainder is
\[
 B(t)=\frac{(t+4)^{3/2}}{3\sqrt t}\operatorname{arsinh}\frac{\sqrt t}{2}
       -\frac49t-\frac43
     =\frac{t^2}{60}+O(t^3).
\]
The resulting weighted propagator is $B(k^2)/(1+k^2)^2$.
The remaining three host loops have eight effective edge parameters.
The sign of a paired term is $(-1)^{|f|}$ for the compatible vertex
forest $f$; the local bubble pair has already been summed.

\paragraph{Overall subtraction and common scale.}
Let $M=\sum_{e=1}^8\alpha_e$, $d$ be the weighted-edge parameter,
$z=F_f/U_f$, $t_0=(M-d)/d$, and $t_1=t_0+z/d$.
Define primitives by
\[
 A'(t)=\frac{B(t)}{(1+t)^2},\qquad
 C'(t)=\frac{tB(t)}{(1+t)^2}.
\]
Then
\begin{align*}
 J(M,d,z)&=d\{t_1[A(t_1)-A(t_0)]-[C(t_1)-C(t_0)]\},\\
 J(M,d,0)&=\partial_zJ(M,d,0)=0,\\
 \partial_z^2J(M,d,z)
 &=\frac{d\,B((M+z-d)/d)}{(M+z)^2},\\
 J(M,d,z)&=d z^2\int_0^1
 \frac{(1-v)B((M+vz-d)/d)}{(M+vz)^2}\,dv .
\end{align*}
This removes the constant and $p^2$ terms of the overall two-point
function. The $p^4$ remainder is retained.

\paragraph{Complete projective measure.}
The full density is
\[
 {\cal D}(\alpha)=\sum_f(-1)^{|f|}
       \frac{J(M,d,F_f/U_f)}{U_f^3}.
\]
Its homogeneity is $1-3(3)=-8$. In each of the eight sectors, set the
largest parameter to one and the other seven to $x_j\in(0,1)$.
Normalization to the simplex gives Jacobian
$(1+\sum_jx_j)^{-8}$, which cancels the density rescaling.
All eight sectors are added once; there is no additional $7!$.
Separate forest terms are not integrated individually.

\paragraph{Independent checks and limits.}
The native implementations use original Gaussian determinants and
positive minor polynomials. An independent reference computes
$F/U$ by solving the Gaussian Schur system and checks the complete
scalar remainder integral. Reusing $A(t_0)$ and $C(t_0)$ within a sector
does not change precision. Original signs, omitted forests and Taylor
orders are deliberately mutated and rejected.
Full-kernel tests and complete-integral acceptance are distinct stages;
the latter still requires the unchanged statistical gates and an
independent stored-data audit.
\end{document}